\documentclass[10pt,a4paper]{amsart}
\usepackage[margin=19mm]{geometry}
\usepackage{amsmath,amssymb,mathtools}
\usepackage[scr=rsfs,cal=euler]{mathalfa}
\usepackage{xfrac}
\usepackage[unicode,hidelinks,bookmarks=false]{hyperref}
\newcommand{\ie}{i.e.\ }
\newcommand{\cf}{cf.\ }
\newcommand{\resp}{respectively\ }
\newcommand{\Subsection}{\subsection}
\providecommand{\nobreakdash}{\nobreak\hbox{-}}
\setlength{\emergencystretch}{2em}
\multlinegap0pt
\usepackage{amssymb}
\usepackage{float}
\usepackage{xcolor}
\usepackage{bm}
\usepackage{enumerate}
\usepackage{algorithm}\usepackage{algpseudocode}
\newtheorem{lem}{Lemma}
\newcommand{\Cay}{{\rm Cay}}
\newcommand{\Irr}{{\rm Irr}}
\newcommand{\Q}{\mathbb{Q}}
\newcommand{\Z}{\mathbb{Z}}
\newcommand{\repF}{\mathfrak{F}}
\title{Character theoretic techniques for nonabelian partial difference sets}
\author{Seth R. Nelson, Eric Swartz}
\date{Source: arXiv:2507.23039v3 (CC BY 4.0)}
\begin{document}
\maketitle

\noindent\emph{Source and licence.} Character theoretic techniques for nonabelian partial difference sets. Version arXiv:2507.23039v3 (CC BY 4.0), distributed by arXiv under the Creative Commons Attribution 4.0 International licence, http://creativecommons.org/licenses/by/4.0/, as recorded in the arXiv licence metadata for this exact version and shown on the abstract page https://arxiv.org/abs/2507.23039v3. Full licence notice: \texttt{licenses/arxiv-2507.23039-CC-BY-4.0.txt}.\par\smallskip

\noindent\textit{Editorial scope.} Counted unit: the article's lem kernel-D (source0.tex lines 364--366). The article's complete original proof of it is delivered with it (source0.tex lines 367--388, one \texttt{proof} environment of 22 lines). No mathematics is written, repaired or re-derived: the statement and the proof are the article's own lines, in the article's own notation, and no retained line is substituted or re-flowed.  Retained uncounted as the source-local context the proof needs: lines 323--325; lines 338--344.  Nothing was omitted from any retained span, so the package carries no omission record.  No retained line contains a citation, so the wrapper carries no bibliography and no bibliography entry.  No retained line contains a figure environment, an image-inclusion command or drawing code, so no image is delivered and the package declares no asset dependency.  Editorial formatting only, in addition: an AMS \texttt{amsart} preamble that restates the article's own declarations and macros used by the retained lines, namely the environments \texttt{lem} on separate counters, together with the standard packages those lines need.  The retained environments print in this document as Lemma 1 in order of appearance, whereas the article prints the same items with the numbering of its own full document.  The complete native source of the article as distributed in the arXiv e-print for this exact version is bundled unchanged as \texttt{sources/182450/source0.tex}.
\begin{lem}\label{phi-function}
	For all $g \in G$,  $\Phi(g) = \sum_{\chi\in\Irr(G)} \overline{\chi(D)} \chi(g)$.
\end{lem}

\begin{lem}\label{eigenvalues}
	For $D$ a regular $(v, k ,\lambda, \mu)$-PDS inside $G$ and a representation $\repF$ affording a nonprincipal irreducible character $\chi$ with underlying vector space $V_\repF$, $\repF(D)$ has two eigenvalues $\theta_1 = \frac12(\lambda-\mu + \sqrt{\Delta})$ and $\theta_2 = \frac12(\lambda-\mu - \sqrt{\Delta})$ with corresponding eigenspaces $V_\repF(\theta_1)$ and $V_\repF(\theta_2)$ such that $V_\repF = V_\repF(\theta_1) \oplus V_\repF(\theta_2)$, and
		\begin{center}
			$\chi(D) = \dim V_\repF(\theta_1)\theta_1 + \dim V_\repF(\theta_2)\theta_2$.
		\end{center}
	In particular, if $\chi$ is a nonprincipal linear character, then $\chi(D) = \theta_1$ or $\chi(D) = \theta_2$.
\end{lem}

	\begin{lem}\label{kernel-D}
		Assume $G$ contains a regular $(v, k, \lambda, \mu)$-PDS with nonprincipal eigenvalues $\theta_1$, $\theta_2$, $\Cay(G, D)$ is not a conference graph, and $\mu > 0$. Let $\mathcal{L}$ be a nontrivial group of linear characters of $G$, and let $\xi\in \mathcal{L}$ be a nonprincipal linear character of order $p$ with kernel $N$. Then, $p | (k-\xi(D))$, $|N\cap D| = \frac{k-\xi(D)}{p} + \xi(D)$, and $|Na \cap D| = \frac{k-\xi(D)}{p}$ for every $a \not \in N$.
	\end{lem}

		\begin{proof}
		Now, $D\not\subseteq\ker\xi$, for $\xi$ is a nonprincipal irreducible linear character. Therefore, $\xi(D) = \theta_i \neq k$ by Lemma \ref{eigenvalues}, so for some $g\in D$ we know that $\xi(g) = \omega$, for $\omega$ a $p^\text{th}$ root of unity. Recall from Lemma \ref{phi-function} that $\overline{\chi(D)} = \sum_{1 = i}^r|h_i^G\cap D|\overline{\chi( h_i)}$ for conjugacy class representatives $h_1, \dots, h_r$. Then,
			\begin{align*}
				0 &= \xi(D) - \xi(D)\\
				  &= \sum_{i=1}^r |h_i^G\cap D|\overline{\xi(h_i)} - \xi(D).\\
				\intertext{We re-express this as a sum over the cosets $Ng_0, \dots, Ng_{p-1}$ for coset representatives $g_0, \dots, g_{p-1} \in G$. In addition, we may order the $g_0, \dots, g_{p-1}$ so that $\overline{\xi(Ng_i)} = \omega^i$.}
				0 &= \sum_{i=0}^{p-1}|Ng_i\cap D|\overline{\xi(Ng_i)} - \xi(D)\\
				  &=\sum_{i=1}^{p-1}|Ng_i\cap D|\omega^i + (|N\cap D| - \xi(D))\omega^0
			\end{align*}
		So, $\omega$ is a root of the polynomial $p(x) = \sum_{i=1}^{p-1}|Ng_i\cap D|x^i + (|N\cap D| - \xi(D))$. Since $\xi(D) = \theta_\alpha \in \Q$ by assumption, then the cyclotomic polynomial $x^{p-1} + \dots + x + 1$ divides $p(x)$, forcing
			\begin{equation*}
				|N\cap D| - \xi(D) = |Ng_1\cap D|  = \dots = |Ng_{p-1}\cap D| = c\in \Z.
			\end{equation*}
			
		Furthermore, 
			
			\begin{equation*}
				k = 1_G(D) = \sum_{j=0}^{p-1}|Ng_j\cap D| = c+\xi(D) + (p-1)c = \xi(D) + pc
			\end{equation*}
			
		so $\frac{k-\xi(D)}{p} = c\in \Z$ and $|N\cap D| = \frac{k-\xi(D)}{p} + \xi(D)$. In particular, $p$ divides $k-\xi(D)$.
		\end{proof}

\end{document}
